\documentclass[11pt]{article}

\usepackage[margin=1in]{geometry}
\usepackage[T1]{fontenc}
\usepackage[utf8]{inputenc}
\usepackage{lmodern}
\usepackage{amsmath,amssymb}
\usepackage{booktabs}
\usepackage{array}
\usepackage{enumitem}
\usepackage{longtable}
\usepackage{xcolor}
\usepackage{hyperref}
\usepackage{listings}

\hypersetup{
    colorlinks=true,
    linkcolor=blue,
    urlcolor=blue
}

\lstset{
    basicstyle=\ttfamily\small,
    breaklines=true,
    columns=fullflexible,
    keepspaces=true,
    frame=single
}

\title{Dynamic Shared Pipeline}
\author{}
\date{}

\begin{document}
\maketitle

\section{What Is Shared}

The dynamic pipeline still prepares one shared branch state and reuses it across all conditions of a dynamic branch execution. That shared state includes:
\begin{itemize}[leftmargin=2em]
    \item the raw traversability matrix,
    \item the static matrix after any preprocessing,
    \item the dynamic loop frames,
    \item the classical and cyclic mapped loops,
    \item one assignment map used for start-goal sampling,
    \item optional spawn masks and campus zone metadata.
\end{itemize}

This is a core scientific contract of the project: dynamic branches compare mappings on the same evolving environment, not on freshly regenerated loops every run.

\section{Preparation Steps}

The current dynamic preparation sequence is:
\begin{enumerate}[leftmargin=2em]
    \item load the source image,
    \item optionally reinterpret the image using campus semantic colors,
    \item optionally preprocess the static obstacle density,
    \item generate the shared dynamic loop,
    \item map that loop into classical and cyclic composite representations,
    \item extract one shared assignment map from the classical loop,
    \item filter allowed spawn vertices and semantic zone vertices on that assignment map.
\end{enumerate}

\section{Cell-Eligibility Modes}

The dynamic pipeline still supports the same high-level eligibility controls:
\begin{itemize}[leftmargin=2em]
    \item \texttt{all\_free}: every thresholded non-black cell is eligible,
    \item \texttt{pure\_white\_only}: only pure-white cells are eligible,
    \item \texttt{zone\_colors\_only}: only designated campus zone cells are eligible.
\end{itemize}

For campus images, white walkways remain traversable but non-spawnable, and gray regions remain non-traversable.

\section{What Changed in the Assignment Layer}

The assignment layer now exposes explicit distribution modes instead of a separate hidden shared-goal toggle:
\begin{itemize}[leftmargin=2em]
    \item start distribution mode,
    \item goal distribution mode,
    \item optional distinct-zone campus sampling,
    \item dispersed sets prioritize internal 8-neighbor separation, with strict enforcement controlled per map type by \texttt{strict\_dispersed\_8\_neighbor\_clearance}; relaxed mode fills any remaining shortage with unique adjacent cells, while clustered sets follow the global \texttt{compact\_clustering} switch,
    \item \texttt{single} target mode for many-to-one shared goals,
    \item one-to-one matching for dispersed and clustered targets, with optional individual reachability.
\end{itemize}

That means the dynamic pipeline still shares \emph{maps}, but each run configuration can now differ in a richer and more explicit way through \texttt{start\_distribution\_mode} and \texttt{goal\_distribution\_mode}.

\section{Solver-Relevant Reuse of Shared State}

The newer low-level guidance logic also benefits from the shared dynamic state. When a dynamic branch enables \texttt{true\_static\_shortest\_path\_distance}, the solver derives a static guidance graph from the already prepared mapped loop and reuses it across repeated low-level replans. When \texttt{tight\_time\_horizon} is enabled, that same static guidance distance is used to build a tighter search horizon for each agent.

So the shared dynamic branch state is not only an experimental fairness device. It is also now the natural source of reusable structural information for low-level dynamic replanning.

\section{Renderer-Only Visual Semantics}

Dynamic campus branches also store \texttt{visually\_free\_vertices}. These are used only by the renderer. A cell can therefore be drawn like an open region even when it is not traversable for pathfinding. This distinction is mainly relevant for campus gray regions.

\end{document}
